\documentclass[11pt]{article}
\usepackage[margin=1in]{geometry}
\usepackage{amsmath,amssymb}
\newcommand{\finalanswer}[1]{\par\medskip\noindent\fbox{\parbox{0.94\linewidth}{\textbf{Answer:} #1}}}
\title{STAT 412 Homework 6 --- Question 19}
\author{}\date{}
\begin{document}\maketitle
\section*{Problem}
For the 10 reserve-capacity values, construct 99\% confidence intervals for $\sigma^2$ and $\sigma$.
\section*{Work}
The sample variance is
\[
s^2\approx0.04636,\qquad df=9.
\]
For 99\%, use $\alpha=0.01$ and the chi-square interval:
\[
\frac{9s^2}{\chi^2_R}<\sigma^2<\frac{9s^2}{\chi^2_L}.
\]
This gives
\[
0.0177<\sigma^2<0.2405,
\]
or rounded,
\[
0.018<\sigma^2<0.240.
\]
Taking square roots gives
\[
0.133<\sigma<0.490.
\]
The correct interpretation for part (a) is choice B: with 99\% confidence, the population variance is between 0.018 and 0.240. Be careful to enter the upper endpoint as 0.240, not 24.

If asked to interpret the standard-deviation interval, use the matching 99\% confidence statement that the population standard deviation is between 0.133 and 0.490 hours.

\finalanswer{Variance CI $(0.018,0.240)$; variance interpretation choice B with between-values $0.018$ and $0.240$; SD CI $(0.133,0.490)$.}
\end{document}
